\section{Related Work}
\label{sec:related}

\paragraph{Math and reasoning retrieval.}
Our target is MathNet-Retrieve~\citep{alshammari2026mathnet}, which asks whether an embedder can match a competition problem to an LLM-written version of the same problem at three levels of disguise. Among neighbouring benchmarks, SABER-Math~\citep{georgiev2026saber} labels human-written problems automatically from start to finish, BRIGHT~\citep{su2025bright} gathers reasoning-heavy queries from real forum posts and curated problem sets, and \citet{wei2026rirsurvey} survey the field. Among retrievers built for reasoning, ReasonIR~\citep{shao2025reasonir} and RaDeR~\citep{das2025rader} train on synthetic queries, the MathLeap embedders come with MELD, a diagnostic for mathematical equivalence~\citep{ye2026meld}, and Rank1~\citep{weller2025rank1} reranks with test-time reasoning; RaDeR, ReasonIR and MathLeap are our $7$--$8$B baselines, none above $0.01$ hard-tier R@1 (Table~\ref{tab:leaderboard-full}).

\paragraph{The blind spot of item-level contamination detection.}
Most contamination research looks for leaked test items. ConStat~\citep{dekoninck2024constat} checks whether a model does better on the benchmark than on reworded copies of it, \citet{yao2024crosslingual} show that a model still remembers a test item after translation, and \citet{yang2023rephrased} show that reworded copies of test items in the training data slip past $n$-gram filters. In all of these the reworded text copies a \emph{benchmark item}, which an item-level detector can catch; our training data rewords public corpus problems. Task contamination~\citep{li2024taskcontamination} and ConStat's synthetic references look past single items to a distribution, and recipe-matching is of that kind, except that the distribution is the benchmark's own writing procedure, which anyone can regenerate rather than leak. What carries over is a style of writing, invisible to item-level checks; leakage explains $1.24$ of the $45$ points (Appendix~\ref{app:contamination}).

\paragraph{Shortcut learning at the procedure level.}
Recipe-matching is a relative of shortcut learning, where a model passes a test by picking up accidental patterns in the data: natural language inference (NLI) models that answer from the hypothesis alone and other crowdworker artefacts~\citep{gururangan2018artifacts,poliak2018hypothesis}, HANS~\citep{mccoy2019hans}, and the framing of \citet{geirhos2020shortcut}. Our style probe is, mechanically, a supervised detector of machine-written text, unlike the zero-shot DetectGPT~\citep{mitchell2023detectgpt}. The usual fix, adversarial filtering, discards the answers or items a shortcut model finds easy~\citep{zellers2019hellaswag,lebras2020aflite}, but it cannot remove a property of how the benchmark was written, which lives in no particular item. Recipe-matching is also not an LLM judge favouring its own outputs~\citep{panickssery2024selfpref,zheng2023judging}, since our generator differs from the benchmark's in vendor and model family, so the model is not recognising itself. What remains is a construct-validity failure~\citep{raji2021everything,bowman2021fixbench}, the benchmark measuring something other than what it claims, here in retrieval evaluation, where BEIR and MTEB~\citep{thakur2021beir,muennighoff2023mteb} set the standard. On MathNet-Retrieve the gaming comes with forgetting of real retrieval, and Appendix~\ref{sec:forgetting} compares two remedies for it, weight interpolation~\citep{wortsman2022wiseft} and replay~\citep{chaudhry2019tiny}.
